\documentclass[11pt]{article}

\usepackage[margin=1in]{geometry}
\usepackage{amsmath,amssymb}
\usepackage{booktabs}
\usepackage{array}
\usepackage{graphicx}
\usepackage{xcolor}
\usepackage[colorlinks=true,linkcolor=blue,citecolor=blue,urlcolor=blue]{hyperref}
\usepackage{natbib}
\usepackage{caption}
\usepackage{enumitem}

\setlength{\parskip}{0.4em}
\setlength{\parindent}{0pt}

\title{\textbf{From-Scratch Small Language Models for Legal and Financial Text:\\
A Staged Pretraining Recipe and an Honest Account of Their Limits}}

\author{Deependra Verma \\
Independent Researcher \\
\texttt{deependra.verma00@gmail.com}}

\date{}

\begin{document}
\maketitle

\begin{abstract}
We build two decoder-only language models from scratch --- 125.8M and 528.5M
parameters --- pretrained on US case law, SEC filings, and educational web
text, then fine-tuned for legal and financial question answering. We report
three findings we believe are broadly useful beyond this specific project.
First, naively mixing five heterogeneous pretraining sources at their natural
size proportions measurably \emph{hurts} held-out quality on the original,
smaller sources (fineweb-edu perplexity worsens from 20.73 to 22.82 despite
more total training data); a staged Warmup-Stable-Decay schedule, adapted
from MiniCPM and SmolLM2, recovers this and improves perplexity on every
shared domain over the smaller model. Second, a subtle supervised
fine-tuning failure mode --- two datasets sharing clause-category names
under different task framings (extraction vs.\ classification) --- causes
the model to collapse to bare label output regardless of what is asked;
we document the bug, its 100\% reproduction rate, and the fix. Third, and
most centrally, we run both models through five externally accepted
benchmark categories via \texttt{lm-evaluation-harness}
\citep{eval-harness} and find a sharp, consistent split: open-book
contract-clause tasks (given a passage, answer about it) are a real,
usable capability, while closed-book legal knowledge and reasoning
(MMLU legal subsets, CaseHOLD) sit at or \emph{below} random chance for
both model sizes. We further report a methodological case study on
legal-domain contamination checking: a naive $n$-gram overlap check flags
an implausible 8{,}712 lines as ``contaminated,'' which on inspection are
uniformly generic legal boilerplate (standard-of-review language, ToS
disclaimer text, canonical precedent quotations) rather than evidence of
benchmark leakage; a stricter 40-word contiguous-run threshold combined
with per-hit provenance attribution narrows this to 5 of 1{,}853 documents
(0.27\%), each independently explainable and none affecting the reported
results. We release both models, the fine-tuning data recipe, an
in-browser ONNX deployment, and all evaluation code.
\end{abstract}

\section{Introduction}

Small (sub-billion-parameter) language models trained from scratch on a
narrow domain are cheap to build and easy to deploy, including entirely
client-side in a browser, but their actual capabilities and failure modes
are not always reported with the same rigor as frontier-model releases.
This paper documents an end-to-end build of two such models for legal and
financial text -- a 125.8M-parameter model and a 528.5M-parameter successor
-- with a specific focus on three questions that matter more for a
deployed small model than raw benchmark leaderboard position: (1) does
adding more pretraining sources actually help, or can naive mixing hurt
it; (2) what training failure modes appear at this scale that would not
be caught by aggregate loss curves alone; and (3) once a model is
fine-tuned into a question-answering assistant, does it hold up against
external, accepted evaluation, not just an in-house test set built by the
same team that built the model.

Our contributions:
\begin{enumerate}[nosep]
  \item Two from-scratch Llama-style decoder language models
    \citep{touvron2023llama} at 125.8M and 528.5M parameters, with a
    shared byte-level BPE tokenizer, trained on US case law, SEC filings,
    and educational web text.
  \item A staged (Warmup-Stable-Decay) pretraining data schedule that
    fixes a directly measured regression from naive proportional mixing
    across five heterogeneous sources, improving held-out perplexity on
    every domain shared between the two model generations.
  \item A supervised fine-tuning recipe combining expert-annotated legal
    datasets (CUAD, LEDGAR) with teacher-distilled question answering
    data, including a documented failure mode -- category-name collision
    between an extraction task and a classification task -- and its fix.
  \item A paired evaluation protocol (both model sizes answering the
    \emph{identical} held-out questions) showing that scaling improves
    aggregate metrics without being uniformly better on every individual
    question.
  \item External validation across five benchmark categories via
    \texttt{lm-evaluation-harness}, with an explicit methodological case
    study on why naive contamination checking over-flags legal text
    specifically, and how to fix it.
  \item Public release of both models, the fine-tuning dataset recipe, an
    ONNX export for in-browser inference, and the full evaluation
    codebase, including the custom benchmark task configurations
    described in Section~\ref{sec:external}.
\end{enumerate}

\section{Related Work}

\textbf{Small and efficient language models.} A line of recent work shows
that careful data curation and training schedules can substitute for
scale: Pythia, OLMo, TinyLlama, Phi, MiniCPM \citep{hu2024minicpm}, and
SmolLM2 \citep{allal2025smollm2} all report staged or curriculum-style
data schedules as part of achieving strong small-model results. We adopt
MiniCPM's Warmup-Stable-Decay (WSD) scheduling idea directly, applying it
not to a learning-rate schedule but to per-source sampling weights, to fix
a data-mixing problem specific to combining a small set of curated,
high-value sources with much larger bulk sources.

\textbf{Domain-specific legal and financial language models.} BloombergGPT
\citep{wu2023bloomberggpt} and SaulLM \citep{colombo2024saullm} both
pretrain at much larger scale (50B and 7B+ parameters, respectively) on
finance- and law-specific corpora. Our work asks a complementary question:
how far can a domain-specific model go at two to three orders of magnitude
fewer parameters, and where precisely does it stop working.

\textbf{Legal NLP benchmarks.} We evaluate against CUAD
\citep{hendrycks2021cuad}, an expert-annotated contract clause extraction
dataset; LexGLUE \citep{chalkidis2021lexglue}, which packages LEDGAR
(clause classification) and CaseHOLD together with several other legal
tasks; the original CaseHOLD dataset \citep{zheng2021casehold}, a
5-way multiple-choice legal holding identification task; and LegalBench
\citep{guha2023legalbench}, a 162-task collaboratively built benchmark
for legal reasoning, from which we construct a 12-task contract-focused
subset (Section~\ref{sec:legalbench}). A concurrent study
\citep{vaddi2026smallmodels} evaluates sub-10B models (including a
3B-active mixture-of-experts model) on ContractNLI, CaseHOLD, and ECtHR
under five prompting strategies including retrieval-augmented generation,
finding that architecture and training quality matter more than raw
parameter count within that range; our work looks two further orders of
magnitude down in scale, at models an order of magnitude smaller than
even that study's smallest model, without retrieval augmentation, to
characterize what remains possible under those constraints.

\textbf{Capacity and memorization.} \citet{allenzhu2024physics} estimate
that language models store roughly 2 bits of knowledge per parameter, a
hard ceiling largely independent of training recipe. \citet{carlini2022quantifying}
show memorization rate scales sharply with how many times a fact is
repeated in training data, with facts seen once memorized at a small
fraction of the rate of facts seen hundreds of times. Both results
directly motivate our closed-book/open-book framing: a model this size
cannot be expected to reliably recall specific, low-repetition legal
facts from parametric memory, regardless of fine-tuning quality.

\textbf{Retrieval-augmented generation.} \citet{lewis2020rag} introduce
retrieval-augmented generation as a way to decouple a model's reasoning
ability from its need to memorize facts. We view retrieval as a
complementary, largely orthogonal direction to the one studied here: this
paper characterizes purely parametric capability at small scale, which we
believe is a necessary baseline before layering retrieval on top, and we
leave a retrieval-augmented extension to future work.

\textbf{General-purpose benchmarks and evaluation infrastructure.} We use
\texttt{lm-evaluation-harness} \citep{eval-harness}, the framework behind
Hugging Face's Open LLM Leaderboard, and its built-in HellaSwag
\citep{zellers2019hellaswag}, ARC-Easy \citep{clark2018arc}, PIQA
\citep{bisk2020piqa}, and MMLU \citep{hendrycks2021mmlu} tasks as a
peer-comparison reference point against other small models, separate from
our legal-domain-specific claims.

\section{Architecture}

Both models share a Llama-style decoder architecture
\citep{touvron2023llama}: RoPE positional encoding \citep{su2021roformer},
SwiGLU activations \citep{shazeer2020glu}, RMSNorm \citep{zhang2019rmsnorm},
and tied input/output embeddings. The 500M model is a straightforward
width-and-depth scale-up; both share the same 16{,}384-token byte-level
BPE tokenizer, trained once and reused rather than retrained, so that
tokenization is not a confound in the size comparison in
Section~\ref{sec:paired}.

\begin{table}[h]
\centering
\caption{Architecture comparison.}
\label{tab:arch}
\begin{tabular}{lrr}
\toprule
 & \textbf{125M model} & \textbf{500M model} \\
\midrule
Parameters & 125,848,320 & 528,538,752 \\
Layers & 12 & 24 \\
Hidden size & 768 & 1{,}152 \\
Attention heads & 12 & 18 \\
Head dimension & 64 & 64 \\
MLP (SwiGLU) intermediate size & 3{,}072 & 4{,}608 \\
Context length (tokens) & 1{,}024 & 1{,}024 \\
Vocabulary & \multicolumn{2}{c}{16{,}384 byte-level BPE (shared)} \\
Positional encoding & \multicolumn{2}{c}{RoPE, $\theta = 10{,}000$} \\
Normalization & \multicolumn{2}{c}{RMSNorm, $\epsilon = 10^{-5}$} \\
Embeddings & \multicolumn{2}{c}{tied input/output} \\
Precision & \multicolumn{2}{c}{bf16 compute, fp32 master weights} \\
\bottomrule
\end{tabular}
\end{table}

\section{Pretraining Data and the Staged Mixing Schedule}
\label{sec:pretrain}

\subsection{Sources}

The 125M model pretrains on three sources: US case law
(\texttt{HFforLegal/case-law}, CC~BY~4.0), SEC 10-K filings
(\texttt{PleIAs/SEC}), and educational web text
(\texttt{HuggingFaceFW/fineweb-edu} \citep{penedo2024fineweb}, ODC-BY), at
a 35\%/42\%/23\% mix, totalling 2.04B unique tokens over 2 epochs
(4.08B tokens seen), reaching a held-out perplexity of 7.76.

The 500M model adds two large bulk sources -- SEC 10-K/10-Q/8-K filings
(\texttt{TeraflopAI/SEC-EDGAR}, Apache-2.0) and additional US case law
(\texttt{common-pile/caselaw\_access\_project}, public domain) -- for
five sources totalling 14.7B unique tokens after cleaning, deduplication,
and 13-gram decontamination against CaseHOLD and LexGLUE (see
Section~\ref{sec:contamination-method} for why this specific decontamination
step matters). Given the larger unique-token budget, this model trains
for a single epoch rather than two.

\subsection{A measured regression from naive mixing}
\label{sec:flatmix}

Our first attempt mixed all five sources proportionally to their natural
size. Because the two new bulk sources are far larger than the three
original sources, this diluted the original case-law/SEC/fineweb-edu
``identity'' of the smaller model down to roughly 13\% of total training
tokens combined. The result was not neutral: held-out perplexity on
fineweb-edu, one of the three original sources, \emph{worsened} from 20.73
(125M model) to 22.82 under this flat mix, despite the 500M model having
more parameters and more total training data. Naive proportional mixing
of heterogeneous sources is not a safe default; it can actively regress
quality on the sources it dilutes.

\subsection{Staged (Warmup-Stable-Decay) fix}

Following the staged-schedule precedent in MiniCPM \citep{hu2024minicpm}
and SmolLM2 \citep{allal2025smollm2}, we replace the flat mix with a
two-phase per-source sampling schedule: broad, size-balanced weights for
the first 85\% of training steps (\emph{stable} phase), then a
\emph{decay} phase for the final 15\% of steps that upweights the three
original, curated sources specifically, so the model sees them last and
retains them most strongly.

\begin{table}[h]
\centering
\caption{Per-source sampling weights, stable vs.\ decay phase.}
\label{tab:wsd}
\begin{tabular}{lrr}
\toprule
\textbf{Source} & \textbf{Stable weight} & \textbf{Decay weight} \\
\midrule
fineweb-edu & 0.15 & 0.30 \\
case-law & 0.13 & 0.25 \\
SEC 10-K & 0.15 & 0.25 \\
SEC-EDGAR (bulk) & 0.30 & 0.12 \\
Case law (bulk) & 0.27 & 0.08 \\
\bottomrule
\end{tabular}
\end{table}

This changes only which windows of already-tokenized data are sampled at
each training step; it required no new data engineering, and did not
change total step count, token budget, or wall-clock time relative to the
flat-mix run.

\subsection{Result: perplexity on every shared domain improves}

\begin{table}[h]
\centering
\caption{Held-out perplexity, 500M (staged schedule) vs.\ 125M.
SEC-EDGAR and case law (bulk) have no 125M-model baseline, since those
sources did not exist in the smaller model's pretraining mix.}
\label{tab:ppl}
\begin{tabular}{lrr}
\toprule
\textbf{Source} & \textbf{125M} & \textbf{500M (staged)} \\
\midrule
US case law & 8.22 & \textbf{7.23} \\
SEC 10-K & 4.36 & \textbf{3.96} \\
Educational web (fineweb-edu) & 20.73 & \textbf{17.81} \\
SEC-EDGAR (bulk) & --- & 3.84 \\
US case law (bulk) & --- & 8.23 \\
\midrule
Aggregate (all sources in each model's val set) & 7.763 & 5.072 \\
\bottomrule
\end{tabular}
\end{table}

The staged schedule not only recovers the flat-mix regression but improves
perplexity over the 125M baseline on all three shared domains. We flag one
honesty caveat explicitly: the aggregate number looks like a larger win
than the per-source numbers alone would suggest, because SEC-EDGAR
(inherently more repetitive, formulaic text) makes up roughly 61\% of the
500M model's validation set by window count. We treat the per-source
numbers, not the aggregate, as the trustworthy signal.

\section{Supervised Fine-Tuning}

\subsection{Data recipe}

We combine two complementary sources of supervision into 21{,}186
instruction pairs (20{,}127 train / 1{,}059 held-out validation):

\begin{itemize}[nosep]
  \item \textbf{Expert-annotated contract data} (ground truth, no
    synthetic generation): CUAD \citep{hendrycks2021cuad} for clause
    extraction, and LEDGAR (via LexGLUE \citep{chalkidis2021lexglue}) for
    clause classification across 100 categories, both CC~BY~4.0.
  \item \textbf{Teacher-distilled question answering}: a locally hosted
    Meta-Llama-3.1-70B-Instruct model generates grounded question-answer
    pairs from passages of the pretraining corpus, judged for quality by
    a combination of NVIDIA Nemotron-Ultra-550B and the same Llama-70B
    teacher, keeping only pairs scoring $\geq 4/5$.
\end{itemize}

Source mix: CUAD ($\sim$51\%), LEDGAR ($\sim$22\%), SEC filings
($\sim$13\%), US case law ($\sim$12\%), educational web ($\sim$2.5\%).
Task mix: extraction ($\sim$60\%), classification ($\sim$22\%), question
answering ($\sim$14\%), summarization ($\sim$5\%), rewriting ($<$1\%).

\subsection{A category-name collision failure mode}
\label{sec:collision}

A handful of clause categories exist under the same name in both CUAD
(where the task is ``extract and quote the relevant text'') and LEDGAR
(where the task is ``name the category'') -- for example \emph{Governing
Law}, \emph{Effective Date}, \emph{Insurance}, and
\emph{Non-Disparagement}. Training on both datasets simultaneously caused
the model to collapse to emitting the bare category name for these
specific clause types \emph{regardless of which task it was actually
asked to perform} -- confirmed with a 100\% reproduction rate on
held-out contracts and hand-written test excerpts before the fix. We
resolved this by removing the roughly 400 conflicting LEDGAR examples for
the affected categories and reformatting all LEDGAR answers into a
structurally distinct template (\texttt{"This clause is categorized as:
X"}) rather than a bare label, so extraction and classification outputs
remain distinguishable throughout training regardless of category-name
overlap. We report this in detail because it is a failure mode that would
not be visible from aggregate training loss and only surfaced through
targeted qualitative inspection of held-out generations.

\subsection{Training procedure}

Both models are fully fine-tuned (not LoRA) from their respective base
checkpoints, at learning rate $3\times10^{-5}$ with cosine decay and 3\%
warmup, bf16 compute with fp32 master weights, loss masked to answer
tokens only. The 500M model's fine-tune ran for 2 epochs; validation loss
ticked up between epoch 1 and epoch 2 (0.1840 $\to$ 0.1887), a directly
measured overfitting signal on a dataset of this size that we take as
evidence against training further epochs without additional data, rather
than as a defect to engineer around.

\section{In-House Evaluation}

\begin{table}[h]
\centering
\caption{In-house held-out evaluation (1{,}059 validation examples),
same methodology for both models: extraction/classification scored
against ground truth, open-ended question answering judged by an
independent local Meta-Llama-3.1-70B-Instruct for semantic correctness.}
\label{tab:inhouse}
\begin{tabular}{lrr}
\toprule
\textbf{Task} & \textbf{125M} & \textbf{500M} \\
\midrule
CUAD clause extraction (Token-F1) & 0.711 & \textbf{0.761} \\
CUAD ``clause not present'' refusal accuracy & \textbf{90.1\%} & 87.7\% \\
LEDGAR clause classification (exact-match) & 74.1\% & \textbf{77.4\%} \\
Case law general QA (closed-book, judged) & 35.6\% & \textbf{42.2\%} \\
SEC filings general QA (closed-book, judged) & 43.1\% & \textbf{46.9\%} \\
Educational web general QA (closed-book, judged) & 18.2\% & \textbf{30.3\%} \\
\bottomrule
\end{tabular}
\end{table}

\subsection{Paired comparison: scaling helps on aggregate, not on every question}
\label{sec:paired}

Aggregate metrics alone can obscure whether a larger model is
\emph{strictly} better or merely better \emph{on net}. We additionally ran
a paired comparison: both models answering the identical 1{,}059 held-out
questions, with both models' answers freshly re-judged in the same pass
to avoid comparing across different judging runs. On CUAD extraction, the
500M model fixed 46 questions the 125M model got wrong, but also newly
got 25 questions wrong that the 125M model had answered correctly --- a
net win, not a clean sweep. One concrete regression, worth citing
directly: asked what changed conditions a defendant cited (extractable
directly from the given source text), the 125M model correctly listed
the real conditions; the 500M model produced a vaguer, legal-sounding but
non-responsive answer. Bigger and better on net does not imply strictly
better on every question, a distinction we believe matters more for a
deployed assistant than the aggregate number alone.

\section{External Benchmark Validation}
\label{sec:external}

The evaluations above are entirely in-house: our own held-out split, our
own extraction/classification scoring, our own judge model. To check the
same claims against outside measurement, we additionally evaluate both
models with \texttt{lm-evaluation-harness} \citep{eval-harness}, the
framework underlying Hugging Face's Open LLM Leaderboard, across five
categories, on the deployed fine-tuned checkpoints.

\subsection{Task selection}

HellaSwag, ARC-Easy, PIQA, and MMLU's \texttt{professional\_law},
\texttt{jurisprudence}, and \texttt{international\_law} subjects are
built into \texttt{lm-evaluation-harness} and run unmodified. CaseHOLD
and LegalBench required custom task configuration, since neither has an
existing harness integration; both configs, along with all raw logs, are
released alongside this paper.

For CaseHOLD, we use its LexGLUE packaging
(\texttt{coastalcph/lex\_glue}, config \texttt{case\_hold}): a 5-way
multiple-choice task where the model must identify the correct masked
holding statement for a case excerpt. We deliberately excluded LexGLUE's
other tasks (SCOTUS, ECtHR, EUR-LEX, UNFAIR-ToS): their documents average
tens of thousands of characters (SCOTUS alone averages roughly 70{,}700
characters, i.e.\ well over 17{,}000 tokens in a sampled check), far
exceeding this model's 1{,}024-token context. Truncating a document to
under 6\% of its length to fit context and then scoring on that truncated
prefix would not produce a number comparable to any published baseline.
CaseHOLD's excerpts, by contrast, average approximately 850 characters
(roughly 200 tokens) and fit comfortably.

\subsection{A curated LegalBench subset}
\label{sec:legalbench}

LegalBench's 162 tasks span formats (free-text generation, multi-label
classification, multiple choice) that do not uniformly fit
\texttt{lm-evaluation-harness}'s simple multiple-choice scoring, and
running the full suite was out of scope for this evaluation pass. We
selected 12 tasks matching two criteria: binary yes/no answer format
(scoreable as simple 2-way multiple choice) and contract-focused subject
matter, matching this project's own CUAD/LEDGAR training domain: nine
ContractNLI clause-entailment tasks, \texttt{contract\_qa},
\texttt{consumer\_contracts\_qa}, and
\texttt{citation\_prediction\_classification}. Each task's prompt is the
benchmark's own official prompt template (including its baked-in few-shot
exemplars), taken verbatim from the LegalBench source repository, not
rephrased, so scores reflect the benchmark's intended evaluation format.

\subsection{Results}

\begin{table}[h]
\centering
\caption{General commonsense reasoning (zero-shot, acc\_norm). Not a legal
claim -- a peer-comparison reference point against other small models.}
\label{tab:general}
\begin{tabular}{lrr}
\toprule
\textbf{Task} & \textbf{125M} & \textbf{500M} \\
\midrule
HellaSwag & 29.4\% & \textbf{32.4\%} \\
ARC-Easy & 41.7\% & \textbf{44.5\%} \\
PIQA & \textbf{60.2\%} & 59.6\% \\
\bottomrule
\end{tabular}
\end{table}

\begin{table}[h]
\centering
\caption{MMLU legal subsets (zero-shot, acc). Random-chance floor for
4-choice is 25\%; both models sit at or near it.}
\label{tab:mmlu}
\begin{tabular}{lrr}
\toprule
\textbf{Subject} & \textbf{125M} & \textbf{500M} \\
\midrule
professional\_law & 24.6\% & 24.7\% \\
jurisprudence & \textbf{30.6\%} & 25.9\% \\
international\_law & \textbf{26.5\%} & 24.0\% \\
\bottomrule
\end{tabular}
\end{table}

\begin{table}[h]
\centering
\caption{CaseHOLD (5-way multiple choice; random-chance floor is 20\%).
Already excluded from both models' pretraining data via 13-gram
decontamination, so no contamination concern applies here.}
\label{tab:casehold}
\begin{tabular}{lrr}
\toprule
\textbf{Metric} & \textbf{125M} & \textbf{500M} \\
\midrule
acc & 11.9\% & 12.1\% \\
acc\_norm & 18.4\% & 19.9\% \\
\bottomrule
\end{tabular}
\end{table}

The 125M model scores \emph{below} the 20\% random-chance floor on
CaseHOLD's acc\_norm metric. Genuinely selecting the correct legal
holding among five plausible-sounding alternatives is a categorically
different, harder skill than the open-book extraction and classification
these models perform well at; we take this as the bluntest possible
external confirmation of a capability boundary this project's own
model cards had already stated qualitatively.

\begin{table}[h]
\centering
\caption{LegalBench, 12-task contract-focused subset (zero-shot,
binary yes/no, 50\% chance floor). Official LegalBench prompts, verbatim.}
\label{tab:legalbench}
\small
\begin{tabular}{lrr}
\toprule
\textbf{Task} & \textbf{125M} & \textbf{500M} \\
\midrule
citation\_prediction\_classification & \textbf{54.6\%} & 50.0\% \\
consumer\_contracts\_qa & 43.2\% & \textbf{58.6\%} \\
contract\_nli\_confidentiality\_of\_agreement & 47.6\% & 50.0\% \\
contract\_nli\_explicit\_identification & \textbf{45.9\%} & 23.9\% \\
contract\_nli\_limited\_use & \textbf{61.5\%} & 60.6\% \\
contract\_nli\_no\_licensing & 59.9\% & \textbf{70.4\%} \\
contract\_nli\_notice\_on\_compelled\_disclosure & 60.6\% & \textbf{66.2\%} \\
contract\_nli\_permissible\_copy & 28.7\% & \textbf{37.9\%} \\
contract\_nli\_return\_of\_confidential\_information & \textbf{54.6\%} & 50.0\% \\
contract\_nli\_sharing\_with\_employees & 50.0\% & \textbf{53.5\%} \\
contract\_nli\_survival\_of\_obligations & 61.2\% & \textbf{71.3\%} \\
contract\_qa & \textbf{48.8\%} & 47.5\% \\
\midrule
\textbf{Mean across 12 tasks} & 51.4\% & \textbf{53.3\%} \\
\bottomrule
\end{tabular}
\end{table}

The LegalBench subset shows a genuinely mixed picture: the 500M model
wins 6 of 12 tasks, the 125M model wins 4, and 2 are effectively tied,
consistent with the paired-comparison finding in
Section~\ref{sec:paired} that scale helps on net without being uniformly
better. One result merits a caveat before being read as a real capability
gap: on \texttt{contract\_nli\_explicit\_identification}, the 500M model
scores 23.9\%, below its own 50\% chance floor, while the 125M model
scores 45.9\%. With $n=109$ examples on a single task, we consider this
plausibly a prompt-format sensitivity rather than a demonstrated
capability regression, and flag it for follow-up rather than drawing a
conclusion from it in isolation.

\section{A Contamination Case Study: Legal Text Needs a Stricter Check}
\label{sec:contamination-method}

Unlike CaseHOLD and LexGLUE, which were excluded from this project's
pretraining corpus by design, LegalBench was never checked for overlap
with the pretraining data before this evaluation. We report the full
investigation because we believe the methodological lesson generalizes
beyond this specific project: naive $n$-gram contamination checks are
too sensitive for legal text specifically, and the standard fix (a
longer $n$) is necessary but not sufficient without also attributing
each surviving hit back to its source.

\subsection{A naive check produces an implausible result}

Following the same method and threshold this project already uses for
CaseHOLD/LexGLUE decontamination (13-word $n$-gram overlap, 64-core
parallel scan across the full $\sim$78GB cleaned pretraining corpus), we
built a contamination set from all 12 selected LegalBench tasks' train
and test splits (1{,}853 documents, 120{,}604 unique 13-grams) and
scanned every source. The result was 8{,}712 overlapping lines across 75
of 76 shard files, spanning every one of the five pretraining sources --
including plain educational web text, which has no obvious reason to
overlap with contract clause data at all. This was too widespread to
trust at face value.

\subsection{Inspection: generic legal boilerplate, not leaked test data}

Rather than accept or reject the result on its aggregate count, we
located and read the actual matched 13-word spans. Every sampled hit was
generic, widely recurring legal language: standard-of-review boilerplate
that essentially every US appellate court recites verbatim in summary
judgment appeals (\emph{``reviewing a motion for summary judgment we
apply the same standard as the [trial court]''}); DMCA/Terms-of-Service
disclaimer boilerplate reused across thousands of unrelated websites; and
canonical precedent-citation language, such as the \emph{Harris
v.\ United States} plain-view-doctrine holding, which any case discussing
that doctrine quotes verbatim. None of these are evidence that a specific
LegalBench document leaked into training -- they are evidence that legal
writing recycles fixed phrases far more than general text does, which a
blunt phrase-overlap check is not equipped to distinguish from genuine
duplication.

\subsection{A stricter threshold, plus provenance attribution}

We re-ran the check requiring a 40-consecutive-word contiguous match --
long enough that coincidental boilerplate reuse cannot plausibly produce
it. This reduced the result to 356 hits, concentrated in a small number
of files rather than spread across nearly all of them. We then attributed
every surviving hit back to its specific source LegalBench document,
rather than treating the count alone as a verdict:

\begin{table}[h]
\centering
\caption{Attribution of the 356 strict (40-word) contamination hits.}
\label{tab:contam}
\begin{tabular}{lrl}
\toprule
\textbf{LegalBench task} & \textbf{Docs affected} & \textbf{What it actually is} \\
\midrule
\texttt{consumer\_contracts\_qa} & 3 / 396 & Facebook's real, public Terms of Service \\
\texttt{contract\_nli\_limited\_use} & 1 / 208 & A standard, widely reused NDA template \\
\texttt{citation\_prediction\_classification} & 1 / 108 & A real historic court opinion, present in \\
 & & both datasets independently \\
\bottomrule
\end{tabular}
\end{table}

In total, 5 of 1{,}853 evaluated documents (0.27\%) show genuine
long-run overlap, concentrated in 3 of the 12 tasks. Each is independently
explainable by the nature of its source: a famous, widely-republished
company Terms of Service document; a standard legal template reused
across many real, unrelated contracts; and a historic court opinion that
both LegalBench (whose \texttt{citation\_prediction\_classification} task
literally consists of real historical case-law excerpts) and our
pretraining corpus (built from the same public case-law record)
independently and unavoidably contain. None of this constitutes narrow
test-set leakage in the sense that matters for benchmark validity, and
none of it is large enough, at under 1\% of any single task's examples,
to move the reported accuracy outside that task's own margin of error.
We consider the LegalBench results in Table~\ref{tab:legalbench} safe to
report on this basis.

We surface this as a general methodological point: legal text's much
higher natural rate of verbatim phrase reuse -- case law quoting
precedent, contracts reusing boilerplate clauses, companies republishing
near-identical Terms of Service -- means a contamination-checking
threshold tuned for general text will systematically over-flag legal
corpora. A longer contiguous-match requirement combined with per-hit
source attribution, rather than raw overlap counts alone, is necessary to
get a trustworthy signal in this domain.

\section{Discussion}

Across both the in-house and external evaluations, a single organizing
finding recurs: these models are a real, usable capability when given a
source passage to work from, and are not a reliable source of legal
knowledge or reasoning from parametric memory alone. This is consistent
with, and adds an external benchmark confirmation to, the capacity-ceiling
argument in \citet{allenzhu2024physics} -- a 528.5M-parameter model can
store at most on the order of 132MB of compressible knowledge, shared
across grammar, style, and every fact it has ever seen -- and the
repetition-dependent memorization curve in \citet{carlini2022quantifying},
which predicts that facts appearing only once or twice in a multi-billion-
token pretraining corpus will be recalled correctly only a small fraction
of the time, exactly the regime most specific legal facts fall into.

We deliberately scope this paper to purely parametric capability, without
retrieval augmentation \citep{lewis2020rag}. This is not a claim that
retrieval would not help -- the open-book/closed-book gap documented
throughout this paper is, if anything, a direct argument \emph{for}
retrieval as a natural next step, since it would convert every closed-book
question into the open-book format these models already handle well. We
leave that extension, and a systematic comparison against it, to future
work, consistent with the recent, larger-scale exploration of
retrieval-augmented small legal models in \citet{vaddi2026smallmodels}.

\section{Limitations}

\begin{itemize}[nosep]
  \item Both models are evaluated against a single held-out validation
    split per model; we do not report cross-validation or multi-seed
    variance, so small differences (a few percentage points) between the
    two model sizes on individual tasks should not be over-interpreted.
  \item The LegalBench subset covers 12 of the benchmark's 162 tasks,
    selected for scoring-format and domain fit rather than
    comprehensiveness; results should not be read as a full LegalBench
    evaluation.
  \item SCOTUS, ECtHR, and EUR-LEX (also part of LexGLUE) were not
    evaluated at all, due to the context-length mismatch described in
    Section~\ref{sec:external}, not because they are uninteresting.
  \item Our contamination check is $n$-gram-based; it would not detect
    paraphrased or semantically equivalent (but lexically different)
    leakage. We consider this an acceptable limitation given the
    combination of a strict contiguous-match threshold and manual
    per-hit provenance review, but note it as a real methodological gap.
  \item Both pretraining runs are single training runs at their
    respective configurations; we do not report variance across
    independent reruns of the same recipe.
  \item Both models are trained on English-language, US-centric legal and
    financial text; none of the findings here should be assumed to
    transfer to other jurisdictions or languages.
\end{itemize}

\section{Conclusion}

We built two from-scratch small language models for legal and financial
text, found and fixed a real data-mixing regression and a real
fine-tuning failure mode along the way, and then subjected the resulting
models to external, accepted benchmark evaluation rather than relying
solely on our own test set. The result is a consistent, externally
validated picture: a genuine, narrow capability at open-book contract
analysis, and an equally genuine absence of reliable closed-book legal
knowledge or reasoning at this parameter count. We release both models,
the fine-tuning recipe, an in-browser deployment, and all evaluation code
so that these specific, quantified claims -- including the ones that do
not flatter the larger model -- can be independently checked rather than
taken on faith.

\section*{Reproducibility}

Model weights, the fine-tuning dataset recipe, ONNX exports for
in-browser inference, and the full pretraining/evaluation codebase
(including every custom benchmark task configuration and contamination-
check script described in this paper) are available at
\url{https://github.com/DeependraVerma/legal-slm-125M}. Model cards with
full evaluation tables are hosted at
\url{https://huggingface.co/DeependraVerma}.

\bibliographystyle{plainnat}
\bibliography{refs}

\end{document}
